%% kapa: exam=1 round=28 year=2017 period=2 form=A subject=appraisal_law subject_ko=감정평가관계법규 seq=19 number=19
%% source: Q-Net 감정평가사 기출문제 (cst00309) artlSeq=5199603 file=extracted/1008/5199603/2017년 제28회 감정평가사 1차시험 문제지/2교시 A형.pdf
%% conversion: kordoc 4.24.51 -> tools/gen1.py (auto); review=proofread-sonnet; flags=none
\begin{question}{19}
부동산 가격공시에 관한 법령상 표준지공시지가에 관한 설명으로 옳지 않은 것은?
\begin{choices}
  \item 표준지에 정착물이 있을 때에는 그 정착물이 존재하지 아니하는 것으로 보고 표준지공시지가를 평가하여야 한다.
  \item 표준지공시지가에 이의가 있는 자는 그 공시일부터 30일 이내에 서면으로 국토교통부장관에게 이의를 신청할 수 있다.
  \item 국토교통부장관은 이의신청 기간이 만료된 날부터 30일 이내에 이의신청을 심사하여 그 결과를 신청인에게 서면으로 통지하여야 한다.
  \item 표준지에 지상권이 설정되어 있을 때에는 그 권리의 가액을 반영하여 표준지공시지가를 평가하여야 한다.
  \item 선정기준일부터 직전 1년간 과태료처분을 3회 이상 받은 감정평가업자는 표준지공시지가 조사ㆍ평가의 의뢰 대상에서 제외된다.
\end{choices}
\end{question}
\answer{4}
